\clearpage
\begin{explanationbox}
  \textbf{\textcolor{CExplanation}{一、为什么会有这条结论？——从折纸开始}}
\begin{center}\begin{tikzpicture}[font=\small,>=stealth,line join=round]
\draw[gray!40,rounded corners=3pt] (0,0) rectangle (14.2,8.3);\node at (7.1,7.700000000000001) {左右等距、函数值同高：对称轴就在中间};
\begin{scope}[shift={(1.7923076923076926,2.038732394366197)},x=1.153846153846154cm,y=0.647887323943662cm]\draw[->,gray] (-0.6,0)--(4.6,0) node[right]{$x$};\draw[->,gray] (0,-0.6)--(0,6.5) node[above]{$y$};\draw[blue,thick,solid,domain=-0.1:4.1,samples=180] plot(\x,{(\x-2)^2+1});\fill[orange!85!black] (0,5) circle(2.2pt);\fill[orange!85!black] (1,2) circle(2.2pt);\fill[orange!85!black] (2,1) circle(2.2pt);\fill[orange!85!black] (3,2) circle(2.2pt);\fill[orange!85!black] (4,5) circle(2.2pt);\draw[violet,dashed] (2,-.5)--(2,6);\node[above,violet] at (2,6){$x=2$};\draw[orange!85!black,dashed] (0,5)--(4,5);\draw[orange!85!black,dashed] (1,2)--(3,2);\foreach \v in {0,1,2,3,4}{\draw[gray] (\v,0)--(\v,-.1);\node[below=5pt] at (\v,-.1){$\v$};}\end{scope}\node[anchor=west,align=left,text width=5.0cm] at (8.4,4.0) {$f=(x-2)^2+1$\\[9pt]$f(0)=f(4)=5$\\[9pt]$f(1)=f(3)=2$\\[9pt]每对横坐标之和为 $4$\\[9pt]轴的位置：$4\div2=2$};
\node[text width=12.6cm,align=center] at (7.1,.65) {图示具体配对；完整证明要求定义域左右成对，且每一对都同高。};
\end{tikzpicture}\end{center}


  把一张纸沿竖直线 $x=2$ 对折：横坐标为 $1$ 的位置会压在 $3$ 上，
  $0$ 会压在 $4$ 上。\textbf{如果图象能够重合，每对位置的高度就必须相同。}
  函数式 $f(x)$ 说的正是“横坐标为 $x$ 时有多高”。
  所以图形上的“等距离、同高度”，翻译成式子就是 $f(2-t)=f(2+t)$。

  \medskip
  
\end{explanationbox}
\clearpage
\begin{explanationbox}
\textbf{\textcolor{CExplanation}{二、先算数，再看图：$f(x)=(x-2)^2+1$}}
\begin{center}\begin{tikzpicture}[font=\small,>=stealth,line join=round]
\draw[gray!40,rounded corners=3pt] (0,0) rectangle (14.2,8.3);\node at (7.1,7.700000000000001) {左右等距、函数值同高：对称轴就在中间};
\begin{scope}[shift={(1.7923076923076926,2.038732394366197)},x=1.153846153846154cm,y=0.647887323943662cm]\draw[->,gray] (-0.6,0)--(4.6,0) node[right]{$x$};\draw[->,gray] (0,-0.6)--(0,6.5) node[above]{$y$};\draw[blue,thick,solid,domain=-0.1:4.1,samples=180] plot(\x,{(\x-2)^2+1});\fill[orange!85!black] (0,5) circle(2.2pt);\fill[orange!85!black] (1,2) circle(2.2pt);\fill[orange!85!black] (2,1) circle(2.2pt);\fill[orange!85!black] (3,2) circle(2.2pt);\fill[orange!85!black] (4,5) circle(2.2pt);\draw[violet,dashed] (2,-.5)--(2,6);\node[above,violet] at (2,6){$x=2$};\draw[orange!85!black,dashed] (0,5)--(4,5);\draw[orange!85!black,dashed] (1,2)--(3,2);\foreach \v in {0,1,2,3,4}{\draw[gray] (\v,0)--(\v,-.1);\node[below=5pt] at (\v,-.1){$\v$};}\end{scope}\node[anchor=west,align=left,text width=5.0cm] at (8.4,4.0) {$f=(x-2)^2+1$\\[9pt]$f(0)=f(4)=5$\\[9pt]$f(1)=f(3)=2$\\[9pt]每对横坐标之和为 $4$\\[9pt]轴的位置：$4\div2=2$};
\node[text width=12.6cm,align=center] at (7.1,.65) {图示具体配对；完整证明要求定义域左右成对，且每一对都同高。};
\end{tikzpicture}\end{center}


  \[
    \begin{array}{c|ccccc}
      x &0&1&2&3&4\\ \hline
      f(x)&5&2&1&2&5
    \end{array}
  \]
  这张表说明：$0$ 和 $4$ 一样高（都是 $5$），$1$ 和 $3$ 一样高（都是 $2$）。
  而且不仅整数如此：
  \[
    f(2+t)=t^2+1,\qquad f(2-t)=(-t)^2+1=t^2+1.
  \]
  不论 $t$ 取哪个实数，这两个数都相等，所以\textbf{整条曲线}关于 $x=2$ 对称。

  \medskip
  
\end{explanationbox}
\clearpage
\begin{explanationbox}
\textbf{\textcolor{CExplanation}{三、图上看：每一对点都在同一水平线上}}
\begin{center}\begin{tikzpicture}[font=\small,>=stealth,line join=round]
\draw[gray!40,rounded corners=3pt] (0,0) rectangle (14.2,8.3);\node at (7.1,7.700000000000001) {左右等距、函数值同高：对称轴就在中间};
\begin{scope}[shift={(1.7923076923076926,2.038732394366197)},x=1.153846153846154cm,y=0.647887323943662cm]\draw[->,gray] (-0.6,0)--(4.6,0) node[right]{$x$};\draw[->,gray] (0,-0.6)--(0,6.5) node[above]{$y$};\draw[blue,thick,solid,domain=-0.1:4.1,samples=180] plot(\x,{(\x-2)^2+1});\fill[orange!85!black] (0,5) circle(2.2pt);\fill[orange!85!black] (1,2) circle(2.2pt);\fill[orange!85!black] (2,1) circle(2.2pt);\fill[orange!85!black] (3,2) circle(2.2pt);\fill[orange!85!black] (4,5) circle(2.2pt);\draw[violet,dashed] (2,-.5)--(2,6);\node[above,violet] at (2,6){$x=2$};\draw[orange!85!black,dashed] (0,5)--(4,5);\draw[orange!85!black,dashed] (1,2)--(3,2);\foreach \v in {0,1,2,3,4}{\draw[gray] (\v,0)--(\v,-.1);\node[below=5pt] at (\v,-.1){$\v$};}\end{scope}\node[anchor=west,align=left,text width=5.0cm] at (8.4,4.0) {$f=(x-2)^2+1$\\[9pt]$f(0)=f(4)=5$\\[9pt]$f(1)=f(3)=2$\\[9pt]每对横坐标之和为 $4$\\[9pt]轴的位置：$4\div2=2$};
\node[text width=12.6cm,align=center] at (7.1,.65) {图示具体配对；完整证明要求定义域左右成对，且每一对都同高。};
\end{tikzpicture}\end{center}


  
  \textbf{看图顺序：}先找中间竖线 $x=2$；再找左右等距的两点；
  最后看这两个点的高度是否一样。图中实线只是帮助看明白，\textbf{不能替代对所有 $t$ 的验证}。
  本包还附有可单独使用的矢量图 \texttt{F020.svg}。

  \medskip
  
\end{explanationbox}
\clearpage
\begin{explanationbox}
\textbf{\textcolor{CExplanation}{四、公式里的 $2a-x$ 从哪儿来？}}
\begin{center}\begin{tikzpicture}[font=\small,>=stealth,line join=round]
\draw[gray!40,rounded corners=3pt] (0,0) rectangle (14.2,8.5);\node at (7.1,7.9) {例6：真实代入位置相加恒为6，中点才是轴};
\draw[->,gray] (1,2)--(13,2) node[right]{$x$};\draw[violet,dashed] (7,1.5)--(7,6.7);\draw[blue,thick] (3,4)--(11,4);\fill[orange!85!black] (3,4) circle(2.4pt);\draw[dashed,gray] (3,2)--(3,4);\node[below=7pt] at (3,2){$t+1$};\fill[orange!85!black] (11,4) circle(2.4pt);\draw[dashed,gray] (11,2)--(11,4);\node[below=7pt] at (11,2){$5-t$};\node[above,violet] at (7,6.7){$x=3$};\node[above=8pt] at (10,4){同高};
\node[text width=12.6cm,align=center] at (7.1,.65) {横向以 $t=0$ 的配对 $1,5$ 示意；一般位置之和恒为6，不指定函数值。};
\end{tikzpicture}\end{center}


  例如对称轴是 $x=3$，左边取到 $x=1$，两点相隔轴的距离都是 $2$，
  右边自然是 $5$。$1+5=2\times3$。
  一般地，已知一边是 $x$，另一边设为 $u$，中间位置应满足
  \[
     \frac{x+u}{2}=a\quad\Longrightarrow\quad u=2a-x.
  \]
  因而 $f(x)=f(2a-x)$ 只是把“取镜子对面的那个位置”写成式子。

  \medskip
  
\end{explanationbox}
\clearpage
\begin{explanationbox}
\textbf{\textcolor{CExplanation}{五、和初中、高中已学知识连起来}}
\begin{center}\begin{tikzpicture}[font=\small,>=stealth,line join=round]
\draw[gray!40,rounded corners=3pt] (0,0) rectangle (14.2,8.5);\node at (7.1,7.9) {平移证明：把原来的对称轴搬到零点};
\node at (3.2,6.9){$f(x)=(x-2)^2+1$};\begin{scope}[shift={(1.2,2.4)},x=1cm,y=.62cm]\draw[->,gray] (-0.5,0)--(4.5,0);\draw[->,gray] (0,0)--(0,6);\draw[blue,thick,domain=0:4,samples=120] plot(\x,{(\x-2)^2+1});\draw[violet,dashed] (2,0)--(2,5.7);\node[below=6pt] at (2,0){$2$};\fill[orange!85!black] (2,1) circle(2pt);\end{scope}\node at (10.3,6.9){$g(t)=f(2+t)=t^2+1$};\begin{scope}[shift={(10.3,2.4)},x=1cm,y=.62cm]\draw[->,gray] (-2.5,0)--(2.5,0);\draw[->,gray] (0,0)--(0,6);\draw[blue,thick,domain=-2:2,samples=120] plot(\x,{(\x)^2+1});\draw[violet,dashed] (0,0)--(0,5.7);\node[below=6pt] at (0,0){$0$};\fill[orange!85!black] (0,1) circle(2pt);\end{scope}\draw[->,orange!85!black,thick] (6.1,4.3)--(7.2,4.3);\node at (7.1,1.5){原图向左平移2个单位；两图横纵单位分别相同，曲线高度不变。};
\node[text width=12.6cm,align=center] at (7.1,.65) {横坐标由 $x$ 改为 $t=x-2$，因此轴 $x=2$ 对应新的 $t=0$。};
\end{tikzpicture}\end{center}


  \begin{itemize}[leftmargin=*,itemsep=0.45em,topsep=0.3em]
    \item \textbf{与二次函数联系：}$(x-2)^2+1$ 的最低点在 $x=2$，左右平方后相同。
    更一般地，$y=Ax^2+Bx+C$（$A\ne0$）可写成
    \[
      y=A\left(x+\frac{B}{2A}\right)^2+C-\frac{B^2}{4A},
    \]
    所以对称轴是 $x=-\dfrac{B}{2A}$。这与本结论是同一个道理。
    \item \textbf{与偶函数联系：}偶函数满足 $g(t)=g(-t)$，对称轴是 $t=0$。
    把原图象的中间位置 $a$ 搬到 $0$，令 $g(t)=f(a+t)$，就得到
    $g(t)=g(-t)$。因此\textbf{关于 $x=a$ 对称，就是“把中间位置移到 $0$ 后变成偶函数”}。
    \item \textbf{与求函数值联系：}若关于 $x=3$ 对称，知道 $f(8)=7$，
    立刻知道 $f(-2)=7$，因为 $8+(-2)=6=2\times3$。
  \end{itemize}

  \medskip
  
\end{explanationbox}
\clearpage
\begin{explanationbox}
\textbf{\textcolor{CExplanation}{六、眼前一亮：两面不同位置的镜子，为什么会产生周期？}}
\begin{center}\begin{tikzpicture}[font=\small,>=stealth,line join=round]
\draw[gray!40,rounded corners=3pt] (0,0) rectangle (14.2,8.3);\node at (7.1,7.700000000000001) {两条对称轴：反射两次，对应一个平移};
\begin{scope}[shift={(2.330769230769231,3.7966666666666664)},x=1.5384615384615383cm,y=1.5333333333333332cm]\draw[->,gray] (-0.8,0)--(3.1,0) node[right]{$x$};\draw[->,gray] (0,-1.4)--(0,1.6) node[above]{$y$};\draw[blue,thick,solid,domain=-0.5:2.8,samples=180] plot(\x,{cos(deg(pi*\x))});\fill[orange!85!black] (0.25,0.7071067811865476) circle(2.2pt);\fill[orange!85!black] (-0.25,0.7071067811865476) circle(2.2pt);\fill[orange!85!black] (2.25,0.7071067811865476) circle(2.2pt);\draw[violet,dashed] (0,-1.3)--(0,1.3);\draw[violet,dashed] (1,-1.3)--(1,1.3);\node[above,violet] at (1,1.3){$x=1$};\end{scope}\node[anchor=west,align=left,text width=5.0cm] at (8.4,4.0) {图取 $f(x)=\cos(\pi x)$\\[9pt]两轴 $x=0,\ x=1$\\[9pt]横坐标配对\\[9pt]$0.25\to-0.25\to2.25$\\[9pt]最后相差 $2$};
\node[text width=12.6cm,align=center] at (7.1,.65) {先关于0反射，再关于1反射；一般得到周期 $2|b-a|$，未必最小。};
\end{tikzpicture}\end{center}


  如果\textbf{同一幅完整图象}既关于 $x=a$ 对称，又关于 $x=b$ 对称（$a\ne b$），
  那么先关于第一条线翻过去，再关于第二条线翻过去，就相当于整幅图向右或向左挪动
  $2|b-a|$。图象仍然重合，因此
  \[
    \boxed{f\bigl(x+2(b-a)\bigr)=f(x)}.
  \]
  也就是说，\textbf{$2|b-a|$ 是一个周期}，不一定是最小正周期。
  例如 $f(x)=\cos(\pi x)$ 的图象关于 $x=0$ 和 $x=1$ 都对称，
  所以 $2$ 是它的周期。

  \medskip
  
\end{explanationbox}
\clearpage
\begin{explanationbox}
\textbf{\textcolor{CExplanation}{七、为什么这里只研究竖直的对称轴？}}
\begin{center}\begin{tikzpicture}[font=\small,>=stealth,line join=round]
\draw[gray!40,rounded corners=3pt] (0,0) rectangle (14.2,8.3);\node at (7.1,7.700000000000001) {水平线对称：普通函数只能留下一个高度};
\begin{scope}[shift={(4.1,2.57)},x=1.0cm,y=0.9199999999999999cm]\draw[->,gray] (-3,0)--(3,0) node[right]{$x$};\draw[->,gray] (0,-1)--(0,4) node[above]{$y$};\draw[blue,thick,solid,domain=-2.5:2.5,samples=180] plot(\x,{2});\fill[orange!85!black] (1,2) circle(2.2pt);\end{scope}\node[anchor=west,align=left,text width=5.0cm] at (8.4,4.0) {若关于 $y=2$ 对称\\[9pt]同一输入 $x$ 的两个高度\\[9pt]$f(x)$ 与 $4-f(x)$ 必须相同\\[9pt]所以 $f(x)=2$};
\node[text width=12.6cm,align=center] at (7.1,.65) {竖轴对称改变横坐标；水平轴对称保留横坐标，同一输入只能有一个高度。};
\end{tikzpicture}\end{center}


  函数 $y=f(x)$ 要求一个 $x$ 只能对应一个 $y$。
  若整幅函数图象关于水平线 $y=b$ 对称，则点 $(x,f(x))$ 的镜像
  $(x,2b-f(x))$ 也得在图上；同一个 $x$ 不允许有两个不同的高度，
  因此必须 $f(x)=b$。换句话说，\textbf{普通函数若关于水平直线对称，图象只能在这条水平线上}。
  这是它和圆、横着开的抛物线等一般曲线的重要区别。
\end{explanationbox}
